\clearpage\thispagestyle{empty}
\pdfbookmark[0]{Sources for the 1805 review of Rohde}{rohde1805-sources}
\begin{center}\textbf{Sources for the 1805 review of Rohde}\end{center}
\noindent This AI translation follows the division title and the review on pp.~483--485 of \href{https://gdz.sub.uni-goettingen.de/id/PPN236007467}{\textit{Werke}, vol.~VI} (Göttingen, 1874). Carl Friedrich Gauss discusses Captain Rohde’s reply to a review of his paper on the masses of the planets. The following collected excerpt from Lindenau’s letter identifies Gauss as the reviewer.
\par\smallskip\noindent The review appeared as part V of XXXIX, \textit{Critik und Anticritik}, in Franz Xaver von Zach’s \textit{Monatliche Correspondenz}, vol.~XII, October 1805, pp.~\href{https://api.digitale-sammlungen.de/iiif/image/v2/bsb10538604_00429/full/1500,/0/default.jpg}{413}--\href{https://api.digitale-sammlungen.de/iiif/image/v2/bsb10538604_00435/full/1500,/0/default.jpg}{419}. The collected text uses dashes where the first printing uses dots.
\par\smallskip\noindent The two French quotations in remark I mean \emph{“to hasten the harvest on all sides, for which one man’s strength is insufficient and mine quite inadequate”} and \emph{“given the large number of new analytical combinations, etc.”} The parenthetical remark between them is Gauss’s. Translation and typesetting produced by AI.
\clearpage
